% !TEX root = ../../tree_notes.tex
% =====================================================================
%  16.2  LCA by Euler tour + RMQ                       sec:c11-euler
%  Exports: def:c11-tour  thm:c11-eulerthm  ex:c11-spq  fig:c11-tour
%           lst:c11-euler  lst:c11-sparse  prob:c11-dist  lst:c11-1135-soln
% =====================================================================
\section{LCA by Euler Tour and Range Minimum}\label{sec:c11-euler}

\begin{previouslybox}[Three questions from 16.1 and earlier chapters]
\begin{enumerate}[nosep, leftmargin=1.8em]
  \item From \(\exref{c11-census}\): \(\lca(8,10) = 1\) because the
        chains \(8\to5\to2\to1\) and \(10\to9\to7\to3\to1\) share only
        the root.  \emph{(If you had a sequence recording every DFS
        visit, where would the shallowest entry between the first
        sightings of~8 and~10 sit?)}
  \item From Ch.~12 (\(\coderef{c09-dfs3}\)): the measuring DFS visits
        each edge twice (down and up) and produces \(2n-1\) node
        recordings in that walk.  \emph{(That is exactly what the tour
        below records.)}
  \item From Ch.~9 (\(\chref{c26}\)): the \emph{Cartesian tree} solves
        RMQ-on-array by building a tree; here we run the reduction in
        the \emph{opposite} direction -- turn LCA-on-tree into RMQ-on-array.
        \emph{(Same algorithm, different port.  Which array do we RMQ?)}
\end{enumerate}
\end{previouslybox}

\begin{topicintro}
The walk of \secref{c11-idea} is correct but \(\bigO{h}\).  This section
converts LCA into a \emph{range-minimum query} on an array derived from
the DFS walk: build the \(2n-1\) tour, mark the first occurrence of each
node, and the LCA is the shallowest tour-entry between the two firsts --
one sentence theorem (\thmref{c11-eulerthm}), then a \(\bigO{n\log n}\)
sparse table answers \emph{any} such query in \(\bigO{1}\) per lookup.
The result: \(Q\) distance/LCA queries in \(O(n\log n + Q)\) total.
Static trees only -- the moment a value changes between queries, hand
off to segment trees (Ch.~19/20).

\textbf{Prerequisites.} \secref{c11-idea} (chains, \(\preceq\), the
census), Ch.~12's DFS (\(\coderef{c09-dfs3}\)) for the walk discipline
and the depth stamp.  Ch.~9's RMQ section is \emph{optional} -- we
rebuild the sparse table here from scratch, but knowing the Cartesian
direction makes the contrast vivid.

\textbf{After this section you can}
\begin{itemize}[nosep]
  \item build the \(2n\!-\!1\) tour, \(\code{first[]}\), and \(\code{dep[]}\)
        by hand on \(T\); \hfill\lv{3}
  \item state and prove the Euler-RMQ reduction (\thmref{c11-eulerthm}),
        and read an LCA off a depth span; \hfill\lv{3}
  \item construct a sparse table in \(O(n\log n)\) and query it in
        \(O(1)\) via two \emph{overlapping} blocks; \hfill\lv{4}
  \item explain why overlapping blocks are legal for idempotent
        operations (\(\min\), \(\max\), \(\gcd\)) but \emph{never} for
        sum, xor, or count; \hfill\lv{4}
  \item solve a multi-query distance problem (CSES~1135) using this
        machine end-to-end. \hfill\lv{3}--\lv{4}
\end{itemize}
\end{topicintro}

% ---------------------------------------------------------------
\subsection{The \(2n-1\) Tour -- What It Is and Is Not}
% ---------------------------------------------------------------

\begin{definition}{Euler tour of a rooted tree}{c11-tour}
Run the root's DFS; append the current node to \(\code{tour}\) \emph{when
you enter it and every time a child's call returns to it}.  For \(n\)
nodes this records \(n\) entries + \(n-1\) returns = \(2n-1\) slots.
Alongside each slot store the \emph{depth} of the node at that slot.
Define \(\code{first}[v]\) = the \emph{earliest} index of \(v\) in
\(\code{tour}\) (0-based).
\end{definition}

\begin{insightbox}[This tour is NOT Ch.~5's flattening tour]
Chapter~5 (\(\chref{c24}\)) records \(n\) entries \emph{only} (enter a
node; never re-log on return) and produces the \(\tin/\tout\) interval
view that makes subtree queries array-range queries.  The \(2n-1\) LCA
tour here is a \emph{different object}: it logs \emph{every} visit (enter
\emph{and} each return), and its whole purpose is to create an array
whose range-minimum of depths answers LCA.  The two tours coexist freely
in the same problem -- \(\tin/\tout\) for ancestry tests and subtree
ranges, the \(2n\!-\!1\) tour for LCA -- but do not confuse them.
\end{insightbox}

% ---------------------------------------------------------------
\subsection{Hand-Computing the Tour on T}
% ---------------------------------------------------------------

The DFS visits children in ascending label order.  Walk \(T\)
(root~1, edges as in \figref{c11-tree}) and record:

\begin{tracebox}[Tour construction on T (0-based index)]
\[
\renewcommand{\arraystretch}{1.15}
\begin{array}{c|ccccccccccccccccccc}
i & 0&1&2&3&4&5&6&7&8&9&10&11&12&13&14&15&16&17&18\\ \hline
\code{tour}[i] & 1&2&4&2&5&8&5&2&1&3&6&3&7&9&10&9&7&3&1\\
\code{dep}[i]  & 0&1&2&1&2&3&2&1&0&1&2&1&2&3&4&3&2&1&0
\end{array}
\]
\medskip
\(\code{first}[\,]\): 1\(\to\)0, 2\(\to\)1, 3\(\to\)9, 4\(\to\)2,
5\(\to\)4, 6\(\to\)10, 7\(\to\)12, 8\(\to\)5, 9\(\to\)13, 10\(\to\)14.

\medskip
\textbf{Verify:} \(\lca(8,10)\): span is \([\code{first}[8],\;\code{first}[10]]
= [5,14]\).  The minimum depth in \(\code{dep}[5..14] = (3,2,1,0,1,2,1,2,3,4)\)
is~\(0\) at \(i=8\), and \(\code{tour}[8] = 1\).  Census says
\(\lca(8,10) = 1\).  \checkmark

\(\lca(4,8)\): span \([2,5]\), \(\code{dep} = (2,1,2,3)\), min~\(1\) at
\(i=3\), \(\code{tour}[3]=2\).  Census: \(\lca(4,5)=2\); since
\(\lca(4,8)=\lca(4,5)=2\) (8 is below 5).  \checkmark

\(\lca(9,10)\): span \([13,14]\), \(\code{dep}=(3,4)\), min~\(3\) at
\(i=13\), \(\code{tour}[13]=9\).  Census: \(9 \preceq 10 \Rightarrow\)
\(\lca=9\).  \checkmark
\end{tracebox}

\begin{figure}[H]
\centering
\begin{tikzpicture}[diagram, scale=0.82, transform shape]
  % depth axis
  \draw[-latex, darkgray] (-0.5,0) -- (10,0) node[right, font=\tiny] {i};
  \draw[-latex, darkgray] (0,-0.5) -- (0,5) node[above, font=\tiny] {depth};
  \foreach \y in {0,1,2,3,4} {
    \draw[gray!25] (0,\y) -- (9.5,\y);
    \node[font=\tiny, anchor=east] at (0,\y) {\y};}
  % tour plotted as a path
  \draw[primaryblue, line width=1pt]
    plot[sharp plot] coordinates{
      (0,0)(1,1)(2,2)(3,1)(4,2)(5,3)(6,2)(7,1)(8,0)
      (9,1)(10,2)(11,1)(12,2)(13,3)(14,4)(15,3)(16,2)(17,1)(18,0)};
  % annotate first[] for 8 and 10
  \fill[accentorange] (5,3) circle (2.5pt) node[above, font=\tiny] {8};
  \fill[accentorange] (14,4) circle (2.5pt) node[above, font=\tiny] {10};
  \fill[successgreen] (8,0) circle (2.5pt) node[below, font=\tiny] {min=1};
  \draw[dashed, warningred] (5,-0.3) -- (5,3.3);
  \draw[dashed, warningred] (14,-0.3) -- (14,4.3);
  \node[font=\tiny, anchor=north, text=warningred!70!black] at (9.5,-0.35)
    {span [first[8], first[10]] = [5, 14]};
\end{tikzpicture}
\caption{The \(2n-1\) depth profile of~\(T\).  The LCA of~8 and~10 is the
node at the shallowest point between their first occurrences: here the
global minimum~0 at \(i=8\), which is node~1.}\label{fig:c11-tour}
\end{figure}

% ---------------------------------------------------------------
\subsection{The Reduction Theorem}
% ---------------------------------------------------------------

\begin{theorem}{Euler-tour LCA = RMQ on depth}{c11-eulerthm}
For any two nodes \(u,v\) in a rooted tree with the \(2n-1\) tour
\(\code{tour}\), depth array \(\code{dep}\), and first-occurrence
\(\code{first}\):
\[
  \lca(u,v) \;=\; \code{tour}\!\bigl[\arg\min_{\code{first}[u]\,\le\, i\,\le\, \code{first}[v]} \;\code{dep}[i]\bigr]
\]
(assuming \(\code{first}[u] \le \code{first}[v]\); swap otherwise).
\end{theorem}

\begin{proofbox}
WLOG \(f_u = \code{first}[u] \le f_v = \code{first}[v]\).  The tour
records the DFS walk.  Between \(f_u\) and \(f_v\) the walk must leave
\(u\)'s subtree (to reach \(v\)) and can only do so by climbing to some
ancestor of~\(u\).  Let \(w\) be the shallowest tour entry in
\([f_u, f_v]\).  Since the walk starts at~\(u\) (of depth \(\depth(u)\)) and
ends at~\(v\), the path \(\code{tour}[f_u..f_v]\) descends into the
subtree(s) and always returns through the same ancestors.  \(w\) is on
the walk from \(u\) up toward the root (so \(w \preceq u\)) and also on
the walk from the root down toward \(v\) (so \(w \preceq v\)); hence \(w\)
is a common ancestor.  If some common ancestor~\(a\) were strictly deeper
than \(w\), both \(u\) and \(v\) would be in \(a\)'s subtree, so the DFS
would visit both without ever surfacing above \(a\) -- contradicting that
\(w\) is the minimum.  Hence \(w = \ell\) (the deepest common ancestor).
\end{proofbox}

% ---------------------------------------------------------------
\subsection{The Sparse Table -- Building and Querying}
% ---------------------------------------------------------------

The theorem reduces LCA to an \emph{RMQ}: given an array \(A[0..N-1]\)
(here \(N = 2n-1\)), answer \(\arg\min\) over any \([L,R]\) in \(O(1)\).
The standard tool is the \term{sparse table}.

\begin{definition}{Sparse table}{c11-spq}
Let \(\code{st}[j][i]\) store the \emph{index of the minimum} \(A\)-value
among \(A[i \;..\; i + 2^j - 1]\) (a \term{sparse-table block} of length \(2^j\)
starting at~\(i\)).  Write \(a = \code{st}[j\!-\!1][i]\) and
\(b = \code{st}[j\!-\!1][i+2^{j-1}]\) (the two halves).  Recurrence:
\[
\code{st}[0][i] = i; \qquad
\code{st}[j][i] = \begin{cases}
  a & \text{if } A[a] \le A[b]\\
  b & \text{otherwise.}
\end{cases}
\]
Build: \(O(N \log N)\), space \(O(N\log N)\).
\end{definition}

\noindent\textbf{Query for \([L,R]\):} let \(j = \lfloor\log_2(R-L+1)\rfloor\).
Two \emph{overlapping} blocks of length \(2^j\) cover \([L,R]\):
\[
  \min = \begin{cases}
    \code{st}[j][L] & \text{if } A[\code{st}[j][L]] \le A[\code{st}[j][R-2^j+1]]\\
    \code{st}[j][R-2^j+1] & \text{otherwise.}
  \end{cases}
\]
Correct because \(\min\) is \emph{idempotent}: duplicating elements in
the overlap does not change the minimum.  \textbf{Warning:} this trick
does \emph{not} work for sum, xor, or count -- their value depends on how
many times an element appears.  Only \(\min/\max/\gcd/\text{and}/\text{or}\)
are safe for overlapping blocks.

\begin{lstlisting}[caption={Build tour, first[] and dep[] on a rooted tree (children ascending)}, label={lst:c11-euler}]
void build_euler(int v, int p, int d) {
    first[v] = (int)tour.size();   // record FIRST occurrence
    tour.push_back(v);
    dep.push_back(d);
    for (int c : adj[v]) {
        if (c == p) continue;
        build_euler(c, v, d + 1);
        tour.push_back(v);         // record RETURN to v
        dep.push_back(d);
    }
}
\end{lstlisting}

\begin{lstlisting}[caption={Build and query a sparse table on the tour depth array}, label={lst:c11-sparse}]
// build: st[j][i] = index of min depth in dep[i .. i+2^j-1]
int lg[N+1];  // precompute lg[1..N]
void build_sparse() {
    int N = (int)tour.size();      // N = 2n-1
    for (int i = 0; i < N; i++) st[0][i] = i;
    for (int j = 1; (1 << j) <= N; j++)
        for (int i = 0; i + (1 << j) <= N; i++) {
            int a = st[j-1][i], b = st[j-1][i + (1 << (j-1))];
            st[j][i] = (dep[a] <= dep[b]) ? a : b;
        }
}
// query: O(1) index of min-depth in [L, R]
int rmq(int L, int R) {
    int j = lg[R - L + 1];
    int a = st[j][L], b = st[j][R - (1 << j) + 1];
    return (dep[a] <= dep[b]) ? a : b;
}
int lca(int u, int v) {
    int L = min(first[u], first[v]);
    int R = max(first[u], first[v]);
    return tour[rmq(L, R)];
}
\end{lstlisting}

\begin{tracebox}[Sparse-table query on T: lca(8, 10)]
\(N = 2n-1 = 19\), span \([L,R] = [5,14]\), length 10.
\(j = \lfloor\log_2 10\rfloor = 3\), block size \(2^3 = 8\).

Left block: \(\code{st}[3][5]\) covers \(\code{dep}[5..12] = (3,2,1,0,1,2,1,2)\),
min at \(i=8\) (depth~0).  Right block: \(\code{st}[3][7]\) covers
\(\code{dep}[7..14] = (1,0,1,2,1,2,3,4)\), min at \(i=8\) (depth~0).
Winner: index~\(8\), \(\code{tour}[8] = 1\).  \(\lca(8,10)=1\).  \checkmark

\(\lca(4,8)\): span \([2,5]\), length 4, \(j=2\), block size~4.
Left \(\code{st}[2][2]\) covers \(\code{dep}[2..5]=(2,1,2,3)\), min at
\(i=3\) (depth~1).  Right \(\code{st}[2][2]\) = same block.  \(\code{tour}[3]=2\).  \checkmark
\end{tracebox}

\begin{complexitybox}[Euler + RMQ method]
\textbf{Setup:} \(O(n)\) tour + \(O(n \log n)\) sparse table on
\(2n\!-\!1\) entries.  \textbf{Query:} \(O(1)\) per LCA.
\textbf{Memory:} \(O(n \log n)\) for the table.
\textbf{Limitation:} static tree only; answers LCA/distance; cannot
produce a \(k\)-th ancestor or a path aggregate.
\end{complexitybox}

% ---------------------------------------------------------------
\subsection{Variations and Twists}
% ---------------------------------------------------------------

\begin{itemize}[nosep]
  \item \textbf{Farmer's linearithmic alternative.}  The \(2n\!-\!1\) tour
        has \(\code{dep}[i+1] - \code{dep}[i] \in \{+1,-1\}\), a
        \emph{\(\pm 1\) RMQ} which admits \(O(n)\) preprocessing and \(O(1)\)
        query (Bender--Farach-Colton).  Contest-irrelevant: the \(O(n\log n)\)
        sparse table already fits in 1~s for \(n\le2\cdot10^{5}\).  \hfill\lv{6}
  \item \textbf{Multiple RMQ on one tour.}  A segment tree over
        \(\code{dep}\) answers \(O(\log n)\) per query but supports point
        \emph{updates} (change a depth = change the tree).  Deferred to
        Ch.~19/22.  \hfill\lv{4}
  \item \textbf{The reverse reduction.}  Ch.~9 (\(\chref{c26}\)) builds a
        Cartesian tree from an array and reads RMQ as LCA \emph{on the
        Cartesian tree}.  That is the mirror image; we do LCA \(\to\) RMQ
        on the \emph{original} tree.  \hfill\lv{3}
\end{itemize}

% ---------------------------------------------------------------
\subsection{Common Pitfalls}
% ---------------------------------------------------------------

\begin{warningbox}[Trap 1: Overlapping-blocks on non-idempotent operations]
If you replace \(\min\) by \emph{sum} in the sparse-table query, the
overlap double-counts.  Only \(\min/\max/\gcd/\text{AND}/\text{OR}\) are
idempotent (\(x \oplus x = x\)); \(\text{SUM}(x,x) = 2x \ne x\).  For
sum/xor on a \emph{static} array, use prefix sums; for path aggregates
use lifting (\secref{c11-aggregates}).
\end{warningbox}

\begin{warningbox}[Trap 2: Forgetting to log the RETURN visit]
The tour must be \(2n-1\) entries, not \(n\).  If you skip the
``push on return'' line, the span between two firsts might \emph{not}
contain the LCA at all -- the DFS went below the LCA and jumped without
logging the climb back.  The inductive invariant: consecutive entries
differ by \(\pm 1\) in depth.
\end{warningbox}

\begin{warningbox}[Trap 3: Using 1-based vs 0-based indexing inconsistently]
The sparse table and \(\code{first}[\,]\) in this section use
\emph{0-based} indices.  Mixing with a 1-based \(\tin\) convention from
Ch.~5 in the same code is a silent off-by-one: \(\code{first}[1]=0\)
(root~1 at tour position~0) vs \(\tin(1)=1\).  Pick one, document it.
\end{warningbox}

% ---------------------------------------------------------------
\subsection{Problem Lab}
% ---------------------------------------------------------------

\solvedproblem{c11-dist}{Distance Queries}{CSES 1135}
  {https://cses.fi/problemset/task/1135}{CSES L3}
  {Euler tour, sparse-table RMQ, distance formula}

\begin{problemstatement}[CSES 1135 Distance Queries]
You are given a tree consisting of \(n\) nodes (rooted at~1) and \(q\)
queries.  Each query consists of two nodes \(a\) and~\(b\), and your task
is to print the distance between them (number of edges on the unique
path).

\textbf{Input.} First line: \(n\) and~\(q\)
($1 \le n \le 10^5$, $1 \le q \le 10^5$).  Next \(n-1\) lines: edges
\(a_i\), \(b_i\).  Next \(q\) lines: queries \(a\), \(b\).

\textbf{Output.} For each query print the distance.

\textbf{Example.}
\begin{verbatim}
Input:
5 3
1 2
1 3
3 4
3 5
1 3
2 5
1 4

Output:
1
3
2
\end{verbatim}
\end{problemstatement}

\begin{attemptgate}
\peek{The formula \(\dist(u,v) = \depth(u) + \depth(v) - 2\depth(\lca(u,v))\)
converts the problem to two depth lookups plus one LCA.}
\nudge{Build the \(2n-1\) tour, sparse table, and answer each LCA in
\(O(1)\) per \thmref{c11-eulerthm}.}
\fullhint{Code: build Euler tour from root~1 via \coderef{c11-euler};
build \coderef{c11-sparse} on the depth array; for each query compute
\(\ell = \lca(a,b)\), print \(\depth(a)+\depth(b)-2\depth(\ell)\).}
\end{attemptgate}

\solstep{Understanding}
The tree is \emph{static} (no updates between queries), so we can
preprocess once.  Each query asks for a path length in edges.  Distance
between two nodes is determined by the hinge (\thmref{c11-eulerthm} +
the formula):
\[
  \dist(a,b) = \depth(a) + \depth(b) - 2\,\depth(\lca(a,b)).
\]

\solstep{Key Observations}
\begin{observation}
The formula is purely \emph{depth}-based.  We need: (i)~the depth of each
node from a single DFS; (ii)~a fast LCA query.  No edge weights, no
aggregates -- the RMQ machine suffices.
\end{observation}
\begin{observation}
Constraints: \(n,q\le 10^5\) rule out \(O(h)\) per query in worst case
(chain \(h = n\)).  The \(O(1)\) sparse-table query is the cleanest fit
because no \(k\)-th-ancestor or path-aggregate is needed.
\end{observation}

\solstep{Think It Through}
\begin{thinkbox}
\deadend{BFS/DFS per query}
{$q = 10^5\), BFS gives \(O(nq) = 10^{10}\).  TLE.}
\deadend{Naive walk}
{Correct, but worst-case \(O(h) = O(n)\) per query on a chain.  TLE.}
\aha{Preprocess the \(2n-1\) tour + sparse table; each LCA is \(O(1)\);
distance is one arithmetic expression.  Total \(O(n\log n + q)\).}
\end{thinkbox}

\solstep{Proof of the Distance Formula}
\begin{proofbox}[Distance through the hinge]
Let \(\ell = \lca(a,b)\).  The route \(a \to b\) goes \(a\to\ell\to b\).
The segment \(a\to\ell\) has \(\depth(a)-\depth(\ell)\) edges (climb the
chain), \(b\to\ell\) similarly \(\depth(b)-\depth(\ell)\).  No
overlap: \(a\)'s segment is below~\(\ell\) and disjoint from \(b\)'s
unless \(a \preceq b\), in which case \(\ell = a\) and the first term is
zero.  Total: \((\depth(a)-\depth(\ell)) + (\depth(b)-\depth(\ell))\).
\end{proofbox}

\solstep{Implementation}
\begin{lstlisting}[caption={CSES 1135 full solution (Euler + RMQ)}, label={lst:c11-1135-soln}]
#include <bits/stdc++.h>
using namespace std;
const int MX = 200001;  // 2n-1 for n=1e5
vector<int> adj[MX], tour, dep;
int firstOcc[MX], st[18][MX], lg[MX], depth_[MX];
int timer_;

void dfs(int v, int p, int d) {
    firstOcc[v] = (int)tour.size();
    tour.push_back(v); dep.push_back(d); depth_[v] = d;
    for (int c : adj[v]) if (c != p) {
        dfs(c, v, d+1);
        tour.push_back(v); dep.push_back(d);
    }
}
void build() {
    int N = (int)tour.size();
    lg[1] = 0;
    for (int i = 2; i <= N; i++) lg[i] = lg[i/2] + 1;
    for (int i = 0; i < N; i++) st[0][i] = i;
    for (int j = 1; (1<<j) <= N; j++)
        for (int i = 0; i + (1<<j) <= N; i++) {
            int a = st[j-1][i], b = st[j-1][i+(1<<(j-1))];
            st[j][i] = (dep[a] <= dep[b]) ? a : b;
        }
}
int lca(int u, int v) {
    int L = min(firstOcc[u], firstOcc[v]);
    int R = max(firstOcc[u], firstOcc[v]);
    int j = lg[R - L + 1];
    int a = st[j][L], b = st[j][R-(1<<j)+1];
    return tour[(dep[a] <= dep[b]) ? a : b];
}
int main(){
    int n, q; scanf("%d%d", &n, &q);
    for (int i = 0; i < n-1; i++) {
        int a,b; scanf("%d%d", &a, &b);
        adj[a].push_back(b); adj[b].push_back(a);
    }
    dfs(1, 0, 0);
    build();
    while (q--) {
        int a,b; scanf("%d%d", &a, &b);
        int l = lca(a,b);
        printf("%d\n", depth_[a] + depth_[b] - 2*depth_[l]);
    }
}
\end{lstlisting}

\solstep{Trace on the Official Sample}
\begin{tracebox}[CSES 1135 sample verification]
Tree: \(1\!-\!2, 1\!-\!3, 3\!-\!4, 3\!-\!5\).  Root~1.
Depths: \(d_1=0, d_2=1, d_3=1, d_4=2, d_5=2\).

Query (1,3): \(\lca(1,3)=1\) (1 is root, hence ancestor of everything).
\(\dist = 0+1-2(0) = 1\).  Expected: 1.  \checkmark

Query (2,5): \(\lca(2,5)\): chains \(2\to1\) and \(5\to3\to1\); intersection
\(\{1\}\); \(\lca=1\).
\(\dist = 1+2-2(0) = 3\).  Expected: 3.  \checkmark

Query (1,4): \(\lca(1,4)=1\).
\(\dist = 0+2-2(0) = 2\).  Expected: 2.  \checkmark
\end{tracebox}

\solstep{Complexity}
\begin{complexitybox}
Tour: \(O(n)\).  Sparse table: \(O(n \log n)\).  Each query: \(O(1)\).
Total: \(O(n\log n + q)\).  Memory: \(O(n\log n)\) for \(\code{st}\).
\textbf{AC.}
\end{complexitybox}

% ---------------------------------------------------------------
\subsection{Practice Corner}
% ---------------------------------------------------------------

\begin{practiceset}
  \practice{Lowest Common Ancestor (many queries)}{SPOJ LCA}
    {https://www.spoj.com/problems/LCA/}{SPOJ L3}
    {the textbook many-query LCA on a judge: build the \(2n\!-\!1\) Euler
     tour and answer each query as a range-minimum-depth over it -- a
     direct transcription of \thmref{c11-eulerthm} with
     \coderef{c11-euler}~+~\coderef{c11-sparse}.  \(N,Q \le 10^{5}\).}
\end{practiceset}

\paragraph{Back-cites.}
Ch.~5 (\(\chref{c24}\)) -- the \(n\)-entry flattening tour and the
\(\tin/\tout\) convention are \emph{different} from the \(2n\!-\!1\) tour
(contrasted above; never re-linked here).
Ch.~9 (\(\chref{c26}\)) -- RMQ via Cartesian tree is the reverse
reduction; we cite the \emph{concept} but own the general sparse table.

% ---------------------------------------------------------------
\begin{recallbox}[Section 16.2 in 10 seconds]
The \(2n\!-\!1\) tour logs every DFS visit (enter \emph{and} return).
\(\code{first}[v]\) is the 0-based index of \(v\)'s first entry.
\thmref{c11-eulerthm}: \(\lca(u,v)\) = tour entry with minimum depth in
\([\code{first}[u], \code{first}[v]]\).  Sparse table answers that in
\(O(1)\) using two overlapping \(2^j\)-blocks (idempotent \(\min\) makes
overlap harmless).  Setup \(O(n\log n)\), query \(O(1)\), static only.
\selftest{
  \item Tour length for \(n = 10\)?  (Answer: \(2(10)-1 = 19\).)
  \item State the LCA theorem in one sentence.
  \item Why are overlapping blocks safe for \(\min\) but not for \(\mathit{sum}\)?
  \item \(\code{first}[8]\) on~\(T\)?  (5.)  \(\code{first}[10]\)?  (14.)
  \item Give the distance formula.
}
\end{recallbox}
